% !TeX root = ../../../main.tex
\section{Methods}

\subsection{Preclinical Model of Ischemic Reperfusion Injury}

Based on Hossain \emph{et al.}~\cite{hossain_Mechanical_2019}, a large animal model of unilateral ischemia-reperfusion injury (IRI) was developed to induce inflammation and/or fibrosis of the renal parenchyma similar to those seen in AMR. Eighteen 3- to 6-month-old Yorkshire cross pigs were purchased from an approved vendor and maintained in a facility accredited by the Association for Assessment and Accreditation of Laboratory Animal Care International. All procedures were approved by the Institutional Animal Care and Use Committee of the University of North Carolina at Chapel Hill. The pigs were fed a standard diet and had access to water \emph{ad libitum}.

For all of the following procedures that used general anesthesia, the pigs were first sedated with Telazol (4.4 mg/kg) and Xylazine (2.2 mg/kg) given together intramuscularly. Subcutaneous atropine (0.04 mg/kg) was administered prior to intubation.  After intubation, anesthesia was maintained with inhaled isoflurane at 2-4\%.

For 15 of the pigs, general anesthesia was induced as described. The femoral artery was isolated by surgical dissection and a 6 to 8 French sheath was placed in the artery. Then, using fluoroscopic guidance, an angioplasty guide catheter was inserted into the sheath and advanced retrograde to the randomly chosen renal artery (left or right). Next, an angioplasty balloon was inserted into the renal artery via the guide catheter and inflated to occlude blood flow for three hours of ischemia. Renal artery occlusion was confirmed by X-ray contrast. Afterwards, the balloon was deflated and removed to allow blood reperfusion. Three additional pigs were not subjected to the balloon occlusion (i.e. no IRI) and were imaged with VisR and DoPIo as controls.

The pigs subjected to IRI were divided into three cohorts: the acute injury, chronic injury, and combined injury (acute and chronic) cohorts. Each group had five pigs with both males and females included. Pigs in the acute injury cohort were allowed to recover for three days and imaged with VisR and DoPIo as described in the following subsection. Since the injured kidneys in this cohort were expected to have undergone acute immune responses, this work refers to these five pigs as the “acute injury cohort”. The next five pigs were assigned to the "chronic injury cohort" where they were allowed to recover from their IRI for 90 days before performing ultrasound elastography studies. The final five pigs, those in the "combined injury cohort", recovered for 90 days after an initial IRI, underwent a second IRI procedure in the same kidneys, and were imaged three days after the second injury so that both acute and chronic injuries were expected to be present. Each cohort was divided so that 2-3 males and female pigs were present in each cohort.

\subsection{Imaging Protocols}

Following the waiting periods described for each cohort, the pigs were imaged with VisR and DoPIo. Each pig was fasted for 12 hours before imaging, and general anesthesia was induced as described above. The anesthetized pigs were then placed in a supine position on an operating table with heating blankets. A flank incision was made on the side of the abdomen with the injured kidney, exposing the organ \emph{in situ}. An ATL L7-4 linear array transducer (Philips Healthcare, Bothell, WA, USA) was placed on the inferior pole of the exposed kidney, and the transducer was oriented to produce a B-mode image of the coronal plane of the kidney. The transducer was positioned such that both the renal cortex and medulla could be visualized at a depth of 25 mm, and that the renal capsule at that depth was tangentially parallel to the vertical axis of the B-mode (i.e. where nephrons were presumably oriented along the lateral axis of the B-mode). Once the initial transducer position was selected, the transducer was affixed to a stereotactic clamp, and the transducer's orientation was calibrated in real-time using an XSens Awinda gyroscope-based motion capture system (Movella Inc., Henderson, NV, USA). The transducer, attached to a Verasonics Vantage 256 scanner (Verasonics Inc., Kirkland, WA, USA) operating the Vantage 4.2.0 software, was then used to acquire B-mode, VisR, and DoPIo images of the kidney in immediate succession using the parameters indicated in Table~\ref{tbl:iriparam}.

\begin{table}[htb]
    \centering
    \caption{Ultrasonic parameters for ARFI excitation, tracking pulses, and displacement tracking.}
    \begin{tabular}{lll}
    \hline
    \textbf{Common Parameters} & Focal Depth (mm) & 25.0 \\
    \textbf{} & Presumed Sound Speed (m/s) & 1540 \\
    \textbf{} & Push Pulse Frequency (MHz) & 4.21 \\
    \textbf{} & Push Pulse Duration ($\lambda$) & 300 \\
    \textbf{} & Push Pulse Voltage (V) & 45.0 \\
    \textbf{} & Push Apodization Function & Hamming \\
    \textbf{} & Track Frequency (MHz) & 6.25 \\
    \textbf{} & Track Pulse Duration ($\lambda$) & 2 \\
    \textbf{} & Track Pulse Voltage (V) & 45.0 \\
    \textbf{} & Track Apodization Function & Hamming \\
    \textbf{} & Track Pulse PRF (kHz) & 10 \\
    \textbf{} & Track Ensemble Duration per A-Line (ms) & 6.0 \\ \hline
    \textbf{VisR-Specific} & Push F-number & 1.5 \\
    \textbf{} & Track Tx F-number & 1.5 \\
    \textbf{} & Track Rx F-number & 1.5 \\
    \textbf{} & NCC Track Kernel Length ($\lambda$) & 3.0 (0.74 mm) \\
    \textbf{} & NCC Search Window Length ($\lambda$) & 0.8 (0.20 mm) \\ \hline
    \textbf{DoPIo-Specific} & Push F-number & 3.0 \\
    \textbf{} & Track Tx F-number & 5.0 \\
    \textbf{} & Track Rx F-number & 1.5, 5.0 \\
    \textbf{} & NCC Track Kernel Length ($\lambda$) & 2.0 (0.49 mm)\\
    \textbf{} & NCC Search Window Length ($\lambda$) & 0.8 (0.20 mm) \\ \hline
    \end{tabular}
    \label{tbl:iriparam}
\end{table}

The initial imaging sequence was repeated three times. The transducer and imaging plane orientations of those acquisitions were denoted as the 0 degree orientation. Using readouts from the gyroscope, the transducer was then rotated about its center axis by 90 degrees while fixed in position to image the kidney in the sagittal plane, and three additional VisR and DoPIo images were taken for each of the aforementioned imaging parameters. The inferior pole of each kidney was marked with a surgical suture at the approximate site of the transducer imaging at the 0 degree orientation.

\subsection{Elastography Data Processing}

Once VisR and DoPIo images were acquired, the data were processed using custom scripts for MATLAB R2024a (MathWorks, Natick, MA, USA). For VisR, the displacement profiles were tracked using normalized cross-correlation (NCC) with parabolic estimation of sub-resolution displacements~\cite{pinton_Rapid_2006} using parameters from previous VisR studies~\cite{hossain_vivo_2017,hossain_Mechanical_2019}; these values are given in Table~\ref{tbl:iriparam}. Parabolic motion filters~\cite{giannantonio_Comparison_2011} were applied using the last 2.5 ms of displacement measurements to remove displacement artifacts from breathing and other external motion, and displacement measurements from each time point were median-filtered with a one-dimensional axial kernel with a 0.5 mm length to enforce smoothness within our measurements' axial resolution. The displacement profiles were then fit to a one-dimensional mass-spring-damper model to calculate relative elasticity (RE), and depth compensation was applied to control for depth-dependent variations in ARF amplitude~\cite{hossain_Evaluating_2018}.

A similar process was applied for DoPIo, noting that displacements were tracked twice due to the use of two distinct tracking focal configurations. For a given ensemble of channel data received from post-ARFI push tracking pulses at a particular beam position, delay-and-sum beamforming was performed twice: once with a narrower (i.e. tightly focused) F/1.5 aperture, and one with a wider F/5 aperture. For each of these paired beamformed A-lines, displacements were estimated by NCC with a kernel size of 2 $\lambda$ (i.e. 2 times the tracking pulse wavelength, as used in previous works~\cite{yokoyama_Assessing_2021, yokoyama_Quantitative_2022,yokoyama_DoubleProfile_2025} and listed in Table~\ref{tbl:iriparam}) and a search window of 0.8 $\lambda$ to account for the lack of a second push as in VisR, then filtered in an otherwise identical manner as VisR. Times-intersect were calculated for the paired displacement datasets, calculated as the time when the two profiles intersected, and the time-intersect was converted into shear elasticity using a linear model derived from finite element model simulations~\cite{palmeri_Ultrasonic_2006} coupled with the Field-II linear ultrasound simulator~\cite{jensen_Calculation_1992, jensen_Field_1997}.

Following the calculation of elasticity metrics, the resulting acquisitions were annotated with 2 mm square regions of interest (ROI) in the outer cortex, inner cortex, outer medulla, and inner medulla. Annotations were made by a single observer who was blinded to the cohort assignments. For each imaging plane, annotations were drawn on a composite of the B-mode and ARFI peak displacement images such that elasticity-related information was made available to the annotator while blinding them from actual VisR and DoPIo data. The ROIs were then applied to VisR and DoPIo images, and the median pixel values of each ROI were extracted. Finally, two additional semi-quantitative metrics were calculated for VisR RE and DoPIo modulus values. First, the degree of anisotropy (DoA) was calculated by computing a ratio of one ROI from a given anatomical region from the 0 degree acquisition versus the same structure from acquisitions at 90 degrees. Furthermore, the regional ratio (RR), where a ratio of one ROI from a given anatomical region versus another ROI from a different anatomical region in the same acquisition (e.g. outer cortex versus inner cortex), was also calculated.

\subsection{Tissue Collection and Histopathology}

Following imaging, pigs in each injury cohort were humanely euthanized with Euthasol while still under general anesthesia. Both kidneys from all three injury cohorts as well as from the uninjured pigs were collected, flushed with 0.9\% saline, then pressure-perfused by injecting 4\% paraformaldehyde for several minutes. Each kidney was stored in 4\% paraformaldehyde for at least 72 hours and then transferred for histology processing. Samples were subsequently placed in 10\% neutral buffered formalin for at least 2 weeks. The posterior and dorsal aspects of the kidneys were marked using suture material at the time of collection.

All tissue trimming, processing, and staining was performed by the Pathology Services Core Facility in the Lineberger Comprehensive Cancer Center at the University of North Carolina at Chapel Hill. Kidney samples for histological evaluation were collected along orthogonal planes. The first section was along the coronal plane at the inferior pole near the sutured site to resemble the imaging plane for the 0 degree transducer orientation. The second section, taken along the sagittal plane, was used to resemble the imaging plane for the 90 degree transducer orientation. Samples were routinely processed to paraffin (ASP6025, Leica Biosystems, Deer Park, IL, USA) and sectioned to 4-5 µm. Samples were histochemically stained using hematoxylin and eosin (H\&E) for routine microscopic evaluation, and picro-Sirius red (counterstained with Weigert's hematoxylin) for collagen (fibrosis) assessment. All slides were scanned digitally using a Aperio AT2 (Leica Biosystems) at 20x magnification.

Immunohistochemistry (IHC) was also performed on the kidney sections to identify mononuclear inflammatory cell infiltrates (i.e. macrophages). IHC slides were dewaxed in Bond Dewax solution (AR9222, Leica Biosystems) and hydrated in Bond Wash solution (AR9590, Leica Biosystems). Heat-induced antigen retrieval was performed at 100\textdegree{}C in Bond-Epitope retrieval solution 1 pH-6.0 (AR9961, Leica Biosystems). After pretreatment, slides were blocked and incubated with Iba1 antibodies (019-19741, Wako Chemicals) at 1:2000 for 1 hour followed by ready-to-use Novolink Polymer (RE7260-K) secondary. Antibody detection with 3,3'-diaminobenzidine (DAB) was performed using the Bond Intense R detection system (DS9263, Leica Biosystems). Stained slides were dehydrated and coverslipped with Cytoseal 60 (23-244256, Fisher Scientific, Waltham, MA, USA). IHC slides, including a positive control slide included for this assay, were digitally scanned in an identical manner as the H\&E and picro-Sirius red slides.

Masked microscopic evaluation was performed by an experienced board-certified veterinary pathologist. H\&E samples were scored qualitatively for tubular changes, inflammation, and glomerular change in the renal cortex and for inflammation and tubular changes in the medulla. Fibrosis was assessed in the cortex and medulla using picro-Sirius red-stained samples. Iba1 positive staining was assessed in both the cortex and medulla, and scoring was based on distribution. All samples were graded using an ordinal scoring system where 0 indicated no findings, 1 indicated minimal findings, and integer rankings increased up to 5 where severe findings were noted. The images were analyzed using QuPath (Queen's University, Belfast, UK)~\cite{bankhead_QuPath_2017} where a board-certified veterinary pathologist assigned ordinal scores to each section based on the degree of inflammation, fibrosis, and macrophage infiltration, and other features observed. The pathologist was blinded to imaging results and cohort assignments. Differences in scoring between the injured kidneys and kidneys from the uninjured cohort were assessed using a Wilcoxon signed-rank test, with a threshold of $\alpha$ = 0.05 to determine significance.

\subsection{Statistical Analyses}

VisR RE and DoPIo modulus estimates, as well as DoA and RR metrics calculated using those parameters, were compared between the injured kidneys of the three IRI pig cohorts and the six total kidneys from the uninjured control group. The likelihood that the distribution of measurements from each cohort came from the identical source was assessed using a Kruskal-Wallis test. Direct comparions between each injured cohort versus the control was performed by applying a Bonferroni post-hoc test to conservatively identify significant differences in those distributions, and indicating comparison pairs where the null hypothesis was rejected given $\alpha$ = 0.05. 

The performance of each classifier was assessed using the area under the receiver operating characteristic (ROC) curve (AUC), and the bias of the model was determined by assessing the specificity and sensitivity at the ROC curve that maximizes Youden's index. The performance of each ROC was assessed through a 10-fold nested cross-validation scheme~\cite{varma_Bias_2006}, leading to a distribution of AUC for each model. This process was repeated for all possible combinations of inputs using raw elasticity metrics, RR, and/or DoA for VisR or DoPIo. Given an input category (e.g. using RR and DoA inputs), the models with the highest mean AUC for each imaging modality (DoPIo v. VisR) were identified, and the significance of difference in their performances were compared by performing a Wilcoxon signed rank test on the AUC distributions, given a threshold of $\alpha$ = 0.05.
